\section{视觉语言模型（VL）：给 AI 装上眼睛}

\subsection{核心目标：Human Level Visual Understanding}

Qwen-VL 的目标是让模型拥有``人类水平的视觉理解''——尽管在某些方面（如细节识别），模型已经超越人类，但在基础的空间关系（上下左右）判断上仍有不足。

\begin{warningbox}{VL 模型的``降质''问题}
长期以来，将视觉能力加入语言模型时会导致语言智力下降（称为``降质''）。Qwen-VL 最新版本首次实现了\textbf{语言理解不降质}——VL 模型的文本能力与 235B 纯文本模型基本持平。
\end{warningbox}

\subsection{四大提升方向}

\begin{enumerate}
    \item \textbf{GUI/手机操控}：提升操控电脑和手机的能力——Agent 需要能``看到''并``操作''屏幕
    \item \textbf{语言智力保持}：VL 模型可以直接当 LM 使用，追赶 Gemini 系列的原生多模态能力
    \item \textbf{多模态 Coding}：代码的输入可以是图像或视频（画个草图就能生成 App），降低 prompt 编写门槛
    \item \textbf{视觉 Reasoning}：结合 Thinking 能力进行视觉推理——例如数一张照片中有多少人（通过逐个标注来计数）
\end{enumerate}

\begin{knowledgebox}{第一人称视角与 VLA}
林俊旸提到一个前沿探索方向：如果有了智能眼镜，每天接收的第一人称视角（Egocentric Video）数据能否用来构建 memory？这驱动了 Qwen 向 VLA（Vision-Language-Action）方向的探索——将机器人控制数据和第一人称视角数据融入模型训练。不过 VLA 目前连基本的 Scaling Law 都未探通。
\end{knowledgebox}

\subsection{本章小结}

Qwen-VL 在保持语言智力不降的前提下，大幅提升了视觉理解能力。下一步是让 VL 模型同时具备视觉 Reasoning 和 GUI 操控能力，并探索 VLA 方向。

